% preamble.tex — Shared LaTeX preamble for Econometrics (Hayashi)

\documentclass[11pt,openright]{book}

%% ── Page Geometry ──────────────────────────────────────────────
% Princeton University Press format, approx 6.125" x 9.25"
\usepackage[
  paperwidth=6.125in,
  paperheight=9.25in,
  top=0.9in,
  bottom=1in,
  left=0.85in,
  right=0.85in,
  headheight=14pt,
  headsep=0.3in,
  footskip=0.4in
]{geometry}

%% ── Encoding & Fonts ──────────────────────────────────────────
\usepackage[T1]{fontenc}
\usepackage[utf8]{inputenc}
\usepackage{mathptmx}           % Times-like serif (matches source)
\usepackage{microtype}

%% ── Math ───────────────────────────────────────────────────────
\usepackage{amsmath}
\usepackage{amssymb}
\usepackage{amsthm}
\usepackage{mathtools}
\usepackage{bm}

% Equation numbering: (chapter.section.number)
\numberwithin{equation}{section}

% Figure/table numbering: chapter.number
\usepackage{chngcntr}
\counterwithin{figure}{chapter}
\counterwithin{table}{chapter}

%% ── Boxed Equations/Definitions ─────────────────────────────────
\usepackage[most]{tcolorbox}
\newtcolorbox{defbox}{
  colback=white,
  colframe=black,
  boxrule=0.5pt,
  arc=0pt,
  left=6pt, right=6pt, top=4pt, bottom=4pt,
  boxsep=2pt
}

%% ── Theorem Environments ─────────────────────────────────────────
\theoremstyle{plain}
\newtheorem{theorem}{Theorem}[section]
\newtheorem{lemma}[theorem]{Lemma}
\newtheorem{corollary}[theorem]{Corollary}
\newtheorem{proposition}[theorem]{Proposition}

\theoremstyle{definition}
\newtheorem{definition}{Definition}[section]
\newtheorem{assumption}{Assumption}[section]
\newtheorem{example}{Example}[section]

\theoremstyle{remark}
\newtheorem{remark}{Remark}[section]

%% ── Problem Set / Exercises ──────────────────────────────────────
\usepackage{enumitem}

% Hayashi uses numbered exercises within "Problem Set" at chapter end
% Format: plain numbers (1, 2, 3, ...) with parts (a), (b), (c)
\newenvironment{problemset}{%
  \section*{Problem Set}%
  \addcontentsline{toc}{section}{Problem Set}%
}{%
}

\newenvironment{answers}{%
  \section*{Answers to Selected Questions}%
  \addcontentsline{toc}{section}{Answers to Selected Questions}%
}{%
}

\newcommand{\exercise}[1][]{\par\medskip\noindent\textbf{#1.}\enspace}
\newcommand{\answer}[1]{\par\medskip\noindent\textbf{#1.}\enspace}

% Starred exercise marker
\newcommand{\starred}{\textsuperscript{*}}

%% ── Figures ──────────────────────────────────────────────────────
\usepackage{graphicx}
\graphicspath{{figures/}}
\usepackage{float}

%% ── Tables ───────────────────────────────────────────────────────
\usepackage{booktabs}
\usepackage{array}
\usepackage{multirow}

%% ── Headers & Footers ────────────────────────────────────────────
\usepackage{fancyhdr}
\pagestyle{fancy}
\fancyhf{}
\fancyhead[LE]{\thepage}
\fancyhead[RE]{\leftmark}
\fancyhead[RO]{\thepage}
\fancyhead[LO]{\rightmark}
\renewcommand{\headrulewidth}{0pt}

\renewcommand{\chaptermark}[1]{\markboth{Chapter \thechapter.\ #1}{}}
\renewcommand{\sectionmark}[1]{\markright{\thesection\ #1}}

\fancypagestyle{plain}{%
  \fancyhf{}%
  \fancyfoot[C]{\thepage}%
  \renewcommand{\headrulewidth}{0pt}%
}

%% ── Cross-references ─────────────────────────────────────────────
\usepackage{hyperref}
\hypersetup{
  colorlinks=false,
  pdfborder={0 0 0},
  pdftitle={Econometrics},
  pdfauthor={Fumio Hayashi}
}
\usepackage{cleveref}

%% ── Index ────────────────────────────────────────────────────────
\usepackage{makeidx}
\makeindex

%% ── TOC depth ────────────────────────────────────────────────────
\setcounter{tocdepth}{2}
\setcounter{secnumdepth}{2}

%% ── Custom Commands (Econometrics notation) ──────────────────────

% Expectation, variance, covariance
\newcommand{\E}{\mathrm{E}}
\newcommand{\Var}{\mathrm{Var}}
\newcommand{\Cov}{\mathrm{Cov}}
\newcommand{\Corr}{\mathrm{Corr}}

% Convergence
\newcommand{\plim}{\operatorname{plim}}
\newcommand{\avar}{\mathrm{Avar}}
\newcommand{\dto}{\overset{d}{\to}}
\newcommand{\pto}{\overset{p}{\to}}
\newcommand{\asto}{\overset{a.s.}{\to}}
\newcommand{\iid}{\overset{\text{i.i.d.}}{\sim}}

% Matrix/linear algebra operators
\newcommand{\tr}{\operatorname{tr}}
\newcommand{\rank}{\operatorname{rank}}
\newcommand{\diag}{\operatorname{diag}}
\newcommand{\vect}{\operatorname{vec}}
\DeclareMathOperator*{\argmax}{arg\,max}
\DeclareMathOperator*{\argmin}{arg\,min}

% Bold vectors
\renewcommand{\vec}[1]{\mathbf{#1}}

% Partial derivatives
\newcommand{\pd}[2]{\frac{\partial #1}{\partial #2}}
\newcommand{\pdd}[2]{\frac{\partial^2 #1}{\partial #2^2}}
\newcommand{\pddm}[3]{\frac{\partial^2 #1}{\partial #2 \partial #3}}

% Ordinary derivatives
\newcommand{\od}[2]{\frac{d #1}{d #2}}
\newcommand{\odd}[2]{\frac{d^2 #1}{d #2^2}}
